\chapter{Evolution of Life}

\section{Natural Selection}

While \textit{Charles Darwin's} theory of evolution has historically been the subject of intense debate, it is no longer a matter of belief. The evidence is overwhelming.
All known life on Earth is based on DNA, and modern biology has uncovered the structure of the genetic code and the mechanisms by which it is replicated, mutated, and selected.
These discoveries strongly confirm the core principles of evolutionary theory.

Complex biological systems, including intelligent and conscious organisms, do not arise spontaneously in their fully developed form.
Instead, they emerge through gradual processes of variation and selection over vast timescales.
The probability of a fully formed human organism appearing by random assembly is effectively negligible.
Rather than appearing suddenly, life evolves through cumulative, incremental changes, where each step is shaped by environmental pressures and reproductive success.


\section{Minimal Conditions for Evolution}

\subsection*{Intelligent Clouds}

There is a science fiction book, \emph{The Black Cloud} \cite{hoyle1957blackcloud},
in which an intelligent cloud arrives and proceeds to cause all sorts
of trouble. The reason the book is science fiction rather than science,
and why intelligent clouds cannot exist in real life, is that there are
no known laws of physics on which such a cloud could plausibly be based.

The more intelligent and complex a system is, the smaller the
probability that it could simply appear spontaneously. For a giant
intelligent cloud, the probability should be practically zero.

The only theory we currently have that explains the existence of
truly intelligent systems is evolution. For evolution to work there
must be many candidates and a mechanism of natural selection capable
of eliminating the less successful ones. In the case of intelligent
clouds, there would have to be lots of clouds, and natural selection
would need to eliminate the weak and disorganized ones while favoring
those capable of maintaining structure. Only then could intelligent
clouds gradually develop.

But how would evolution work for a cloud whose behavior is governed
by the Navier--Stokes equations of fluid dynamics? How could such a
cloud keep its information in order? What would happen if a storm
passed by and blew the cloud apart? Even relatively small disturbances to a human brain can cause severe
damage. A violent disturbance would correspond to putting a brain into
a kitchen blender and switching it on. It is not difficult to see why
the brain would not think very clearly afterward.
An intelligent cloud would need to know its boundaries and maintain
a stable internal structure to stay 'alive'. A gas cloud has neither. It would not
remain organized long enough to evolve intelligence.

In principle one could imagine some unknown form of matter with
properties suitable for maintaining information in a gaseous state.
However, astronomical observations give us little support for this
idea. When we analyze the spectrum of light coming from anywhere in
the universe—even from the most distant galaxies—we see essentially
the same electromagnetic fingerprints. This tells us that they are
made of the same raw materials as our own home sweet Milky Way.



\subsection*{Alien Life}

It has been proposed that life could be far more common than generally believed and might not necessarily be carbon-based.
However, the challenge is that Earth appears to be an ideal habitat for diverse life forms.
If life emerged easily, shouldn't we observe "exotic" organisms—perhaps non-DNA-based—coexisting with us? Instead, we find only one biochemical lineage, all rooted in DNA and evolution.
Furthermore, our exploration of the solar system and our extensive monitoring of radio signals (SETI) have yielded no traces of life, intelligent or otherwise.

Perhaps life exists on a different temporal scale. If their biological processes run significantly faster or slower than ours, we might perceive them as either stationary matter or a blur too rapid to recognize as sentient.
However, while time-scale differences are a great sci-fi concept, biology is bound by the laws of thermodynamics and chemistry.
Chemical reactions (the basis of life) happen at specific rates dictated by temperature and molecular stability.
A creature "moving too fast" would likely burn up from the heat of its own metabolism; one "moving too slow" might not be able to gather enough energy to maintain its structure against entropy.

Maybe exotic life could be here, but we simply aren't looking for it correctly.
Most of our tools (PCR, DNA sequencing) are designed specifically to find DNA.
If a "non-DNA" microbe existed in the dirt, our current tests would likely dismiss it as "non-living" chemical noise.


\subsection*{Super-Symmetric Life}

The vast majority of matter in the universe consists of "dark matter," the nature of which remains one of science's greatest mysteries.
Some theories of physics predict (or assume) a set of supersymmetric particles that could account for this phenomenon.
If these particles are capable of forming complex structures—analogous to atoms and molecules—then it is statistically plausible (given that dark matter is five times more prevalent than visible matter) that "dark life" exists.
We may be sharing the universe with an entire dark ecology that remains completely invisible to our senses.

Again, physics fights back. Dark matter (and most predicted supersymmetric particles) can be shown to be collisionless.
It doesn't interact with electromagnetism, which means it doesn't "clump" the way normal matter does.

To have life, one needs complex molecules. Normal matter forms molecules because electrons attract and repel each other.
Since dark matter doesn't seem to interact via the electromagnetic force, it can't form "dark atoms" or "dark DNA."
It mostly just passes through itself and us like a ghost. Without a way to bond particles together, dark matter might not be able to build a "dark person."


\section{Definition of Living Observer}

Defining life itself is notoriously difficult.

The primary challenge stems from borderline cases.
Entities such as viruses, prions, and sterile organisms satisfy some criteria for life but fail to meet others.
Any strict definition inevitably excludes entities that many scientists consider "living," or includes those they do not.

Recent developments in AI and computing present a distinct challenge: substrate independence.
We can create digital replicas of living cells; in the future, we may run increasingly complex simulations of life.
These simulated entities do not truly exist, appearing only as mathematical equations executed by silicon hardware.

Do these digital replicas qualify as life?

Some non-living systems—such as growing crystals, spreading fire, or autocatalytic chemical sets—exhibit superficial signs of life.
However, they lack the active, self-directed control that characterizes true living systems.
Passive or fleeting separation from their environment is insufficient.

Despite these complexities, certain characteristics appear consistently across all known living systems:

\begin{itemize}
\item The Cellular Unit: All life on Earth is founded upon cellular structures.

\item Thermodynamics: The second law of thermodynamics dictates that complex systems naturally trend toward disorder.
  By favoring structures that efficiently preserve and replicate information, evolutionary processes gradually accumulate complexity.
  DNA appears to play a critical role here. Life operates as a low-entropy system, continuously exporting disorder to maintain internal order.

\item Intelligence: Some degree of intelligence, or at minimum, automated self-regulation, is required.
  An extreme example of this is the conscious, intelligent entity that attempts to describe its own existence mathematically through the lens of Quantum Mechanics—the observer. 

\item Informational Boundaries: Life requires well-defined boundaries between the "inside" and the "outside."
  In fact, even minor breaches are often lethal. For example, if human skin is wounded, the consequences can be catastrophic.
  The constituent atoms remain, but the intricate internal organization begins to dissipate, leading to systemic failure.
\end{itemize}


\section{Conclusion}

\begin{principle}{Observer Requires Informational Boundaries}
{observer-boundaries}

An observer requires informational boundaries to preserve coherent internal state.

\end{principle}
